% Background and Motivation: two background topics, two motivations,
% followed by a standalone failure-position example introducing AegisKV.
\section{Background and Motivation}
\label{sec:background-motivation}

\subsection{Background}
\label{sec:motivation-background}

\paragraph{Multi-agent workflows.}
Multi-agent workflows are becoming an important form of LLM application,
including coding assistants and research systems~\cite{wu2023autogen,anthropic2025multiagent}.
An agent may have a predefined role, such as a planner or reviewer, or be
instantiated as a subagent to complete a particular subtask. Each agent
may issue several LLM requests, invoke tools, and incorporate their results
before producing its output. These interactions form a workflow with
sequential dependencies, parallel branches, and joins. Multiple requests
using the same role may share a compatible prompt prefix, while successive
tool interactions extend an individual agent's history. Thus, a request,
an agent, and the complete user workflow describe different units of
execution and completion.

\paragraph{Failures at serving scale.}
LLM services distribute requests across many GPU workers, with each model
replica potentially spanning several devices. As deployments grow,
failures become an operational concern even when individual devices are
reliable: an unavailable device can interrupt an entire replica and shift
its work to the surviving fleet~\cite{cao2026lumen}. The interruption loses
both serving capacity and access to the failed worker's GPU-resident KV
cache. Recovery therefore requires more than rerouting a request; it must
also supply the state needed to resume computation. KV replication and
checkpointing systems already address this state-recovery
problem~\cite{strati2024dejavu,jayakody2026ghostserve,cao2026lumen}.
We focus on fail-stop GPU/worker failures with a compatible surviving
replica, retained token history and control state, and surviving backup
storage. Under this scope, the application-level consequences of the
interruption depend on both its workflow position and its recovery cost.

\begin{figure*}[t]
  \centering
  \captionsetup[subfigure]{font=normalsize,labelfont=normalfont,
    textfont=normalfont,justification=centering,singlelinecheck=false}
  \begin{subfigure}[t]{0.32\textwidth}
    \vspace{0pt}
    \begin{minipage}[c][2.25in][c]{\linewidth}
      \centering
      \includegraphics[width=\linewidth]{figures/workflow.pdf}
    \end{minipage}
    \caption{Four-branch workflow}
    \label{fig:motivation-workflow}
  \end{subfigure}\hfill
  \begin{subfigure}[t]{0.32\textwidth}
    \vspace{0pt}
    \begin{minipage}[c][2.25in][c]{\linewidth}
      \centering
      \includegraphics[width=\linewidth]{figures/workflow-gaps.pdf}
    \end{minipage}
    \caption{Workflow-maximum gaps}
    \label{fig:motivation-branch-gap}
  \end{subfigure}\hfill
  \begin{subfigure}[t]{0.32\textwidth}
    \vspace{0pt}
    \begin{minipage}[c][2.25in][c]{\linewidth}
      \centering
      \includegraphics[width=\linewidth]{figures/recovery-cost.pdf}
    \end{minipage}
    \caption{Context and Prefill cost}
    \label{fig:motivation-recompute}
  \end{subfigure}
  \caption{Motivating workflow-aware KV protection.
    (a) Selected Claude Code trace example with four parallel branches;
    A0/A5 denote Main before/after them. Bars preserve observed timing;
    dashed tails assume that all branch outputs are required at the join.
    (b) ECDF of the maximum branch-duration gap in each of 121 workflows
    with temporally inferred fan-in; it is not a slack distribution.
    (c) Measured Prefill latency for Qwen2.5-7B-Instruct-1M on two
    A100 PCIe 40\,GB GPUs (TP2, BF16, batch one), with MInference/DCA
    sparse Prefill above 32K context, including the final one-token EXTEND.
    Reuse starts from the same $(C-1)$-token KV and computes
    one new token. DRAM restores all historical KV by pinned H2D,
    without transfer--compute overlap.
    Full prefill starts from empty KV with 8K chunks and the same final
    one-token boundary. Reuse points are medians and min--max ranges of
    three timings; full prefill has one timing per length.
    K~=~1,024; M~=~$10^6$.}

  \label{fig:motivation}
  \Description{Left: a four-branch workflow with A2 completing last at
    151.6 seconds. Middle: an ECDF of workflow-maximum branch-duration
    differences with a median of 138 seconds and a 90th percentile of
    1173 seconds. Right: three measured Prefill curves for Qwen2.5-7B-Instruct-1M
    from 8K to 1,010,000 tokens. Blue circles denote HBM reuse,
    orange squares DRAM reuse without overlap, and red triangles
    full prefill with sparse attention above 32K. Both axes are logarithmic.}
\end{figure*}

\subsection{Motivation}

\paragraph{Workflow completion requires different protection priorities.}
\label{sec:motivation-workflow}
The user waits for the complete task, so we use workflow completion time
(WCT), from task release to completion of all required results, as the
primary objective. Minimizing an individual request's or agent's completion
time alone does not capture the effect of its dependencies.
Figure~\ref{fig:motivation}(\subref{fig:motivation-workflow}) illustrates
this distinction with a four-branch coding workflow. A2 completes last,
151.6\,s after the fork, whereas A1 completes at 60.4\,s. If the following
Main stage needs all four outputs, accelerating A1 alone does not advance
the join, while delaying A2 postpones downstream work. The same distinction
matters for fault protection: branches differ in how much additional delay
they can absorb. We call this tolerance \emph{slack}.

To characterize this heterogeneity beyond one example, we analyze the
public SemiAnalysis CC/Weka dataset collected from Claude Code coding
workflows~\cite{semianalysis2026cctraces}. Its filtered release contains
393 sessions. We retain overlapping subagent groups whose branches
complete and are followed by a Main request after the last branch finishes;
Main may also execute while the branches run. This yields 286 groups in
121 workflows. For each group, we compute the longest minus the shortest
branch duration; each workflow contributes its maximum such difference.
Figure~\ref{fig:motivation}(\subref{fig:motivation-branch-gap}) shows a
median of 138\,s and a 90th percentile of 1,173\,s. These statistics show
branch-duration heterogeneity among workflow maxima. The trace lacks
explicit result-consumption edges, so fan-in is inferred from timing and
the gaps are not a distribution of slack at failure time. The example and
distribution together motivate assigning protection priorities according
to workflow impact, rather than treating all agent requests equally.

\paragraph{Growing contexts make KV reconstruction increasingly costly.}
\label{sec:motivation-long-context}
An agent's later requests often include earlier reasoning, code, tool
outputs, and other agents' results. Without context reduction, this history
grows across interactions. KV caching avoids repeatedly processing that
history during normal execution. If a failure makes the active KV
unavailable and no usable copy survives, recovery must rebuild the missing
state before continuing; longer histories can therefore turn a brief
interruption into substantial additional computation.

To isolate the cost of populating KV,
Figure~\ref{fig:motivation}(\subref{fig:motivation-recompute}) compares
three paths: HBM reuse, DRAM reuse without transfer--compute overlap,
and full prefill from empty KV. We measure
Qwen2.5-7B-Instruct-1M~\cite{qwen2025qwen25oneM} on two A100 PCIe
40\,GB GPUs with TP2, BF16 weights and KV, and batch one. The SGLang
0.5.10rc0 worker uses MInference's vertical/slash sparse
attention~\cite{jiang2024minference} with the published Qwen1M head
configuration and DCA position handling. We retain the 32K context
threshold; above it, even the final one-token EXTEND uses sparse attention.
At or below the threshold, all paths use grouped dense DCA.
The sparse implementation shares GQA K/V without materializing repeated
heads. Prefix chunks group at most 2 query heads to bound selector
memory; the final query batches all heads. All three paths use the same
implementation and configuration.

Total context is $C=M+1$. Full prefill processes $M$ tokens in chunks of
at most 8,192, then executes one separate token. Each reuse path starts
from that same sparse-generated historical KV and executes the same
one-token EXTEND. HBM reuse has the prefix resident before timing.
DRAM reuse copies the complete historical KV from pinned host memory
and synchronizes the copy stream before computation. It neither
compresses KV nor loads only selected sparse blocks.
Every path ends at synchronized model-body completion with all $C$
tokens' KV resident, before vocabulary projection. Model loading and
backup creation are excluded; no failure is injected.

At 1,010,000 tokens, HBM reuse, DRAM reuse, and sparse full prefill
take 1.266\,s, 2.517\,s, and 565.93\,s, respectively. KV reuse avoids
historical computation even with sparse Prefill enabled. HBM reduces
latency by 49.7\% relative to serialized DRAM in this configuration.
This comparison isolates a non-overlapping restore path; it does not
quantify HBM's advantage over a pipelined DRAM implementation.

The figure uses one fixed synthetic text stream. Each reuse point is
the median of three timings after one warmup; error bars show min--max
ranges, not independent-workload uncertainty. Full prefill has one
timing per length. Short forwards and operator checks precede timing,
including a 1M-context, 8K-query operator check. A complete original
all-head sparse model reference additionally precedes the 64K point;
no separate full-prefill warmup is run at 256K or 1.01M.
We check matched final hidden states, all-layer new-token KV, and greedy
outputs, and verify historical KV bitwise after restoration. CUDA events
confirm that DRAM transfer completes before computation. These
checks validate the measured paths under the same sparse configuration;
they do not establish equivalence to dense attention or long-context
task quality. Complete-history HBM is a reference state, not a
finite-HBM protection policy, and worker-level latency does not establish
end-to-end WCT benefits.

\begin{figure}[t]
  \centering
  \includegraphics[width=\linewidth]{figures/recovery-exposure.pdf}
  \caption{Workflow position changes the impact of recovery.
    Four independent single-branch failure counterfactuals for
    Figure~\ref{fig:motivation}(\subref{fig:motivation-workflow}),
    at 29.5\,s after the fork. Hatching is delay absorbed by join slack;
    solid fill is WCT loss. Reconstruction uses the active request's
    64-token hash-block input proxy and an assumed 64-token tail, with
    Qwen3-14B/A100 BF16 static estimates (batch and TP one).
    Other execution and queueing are held fixed; this is not a
    simultaneous four-branch fault experiment.}
  \label{fig:motivation-recovery-exposure}
  \Description{For A1 through A4, reconstruction costs are 35.2, 32.2,
    36.7, and 38.9 seconds, while corresponding WCT losses are 0, 32.2,
    17.6, and 17.0 seconds. A1 absorbs all reconstruction delay;
    A2 has zero slack and exposes all of its delay.}
\end{figure}

\paragraph{Failure position determines the workflow benefit of protection.}
The two motivations must be considered together: expensive reconstruction
does not imply an equally large WCT penalty. Figure~\ref{fig:motivation-recovery-exposure}
places an independent hypothetical failure on each active branch of
Figure~\ref{fig:motivation}(\subref{fig:motivation-workflow}) at the same
29.5\,s point. Under the stated join assumption, the observed completion
times provide offline slack, and static estimates provide the additional
reconstruction delay. A1's 35.2\,s cost fits within its 91.2\,s slack,
so it adds no WCT in this counterfactual. A2 costs only 32.2\,s to rebuild
but has zero slack, exposing the entire delay. A3 and A4 expose 17.6 and
17.0\,s, respectively. Hence, the branch with the cheapest reconstruction
can be the most important to protect. This controlled example illustrates
dependency effects; concurrent failures and shared recovery resources can
also delay other requests and must be evaluated separately.

\label{sec:motivation-goals}
These observations lead to AegisKV: a workflow-aware KV protection system
that allocates limited recovery state according to its effect on WCT.
The design connects request-level delay tolerance to the recovery value
of that request's KV, while accounting for the HBM, DRAM, and bandwidth
consumed by protection. Existing valid cache copies remain usable for
recovery; additional protection must justify its cost to normal serving.
The next section develops the corresponding protection, placement, and
recovery mechanisms.
